\section{A genuine action polar need not have a continuous-form primitive}
\label{sec:r24-divisor-polar-primitive-obstruction}

\paragraph{Status and precise failed inference.}
This is a \texttt{FAILED\_ATTEMPT} audit of the automatic-existence
step one might try to add to VP. Automatic supply of the specified
continuous-coefficient primitive is \texttt{FALSE}. We first prove
this for the explicit canonical representing current; the final
exact Hessian probes also exclude any alternative representative
satisfying both VP hypotheses. No uniqueness of current
representatives is assumed. The sufficient implication in VP
is unchanged.
No actual analytic relation with positive carrier seminorm is
constructed, and KC12, CM10 and CTO remain \texttt{OPEN}. This
companion has zero principal-theorem and zero result-family count.
All spaces, currents, smooth tests and compact tubes are over the
original \(K/\mathbf Q_p\) with its fixed absolute value.

Put
\[
 Y=(\mathbf G_{m,K}^2)^{\mathrm{an}},\qquad
 X=Y\times_K\mathbf D,\qquad
 i(y)=(y,0),\qquad p:X\longrightarrow Y,
 \tag{VE24.1}
\]
where \(\mathbf D=\{|T|<1\}\) is the original open unit disc.
Write \(s_j=\log|Z_j|\) for the two coordinates of \(Y\).
The analytic product, rather than a set-theoretic product of points,
is used throughout. The integration current
\(Q=i_*[Y]\) is an allowed original symmetric middle-positive,
separately closed current of degree \((1,1)\) on the pure
three-dimensional space \(X\). Its exact support is
\(F=i(Y)\). For the polar normalization below we use the
whole-domain absolute carrier \(G=X=\operatorname{supp}[X]\).
Thus no identification of the analytic section with a unit-hinge
absolute cap is required; \(Q\) is carried by the allowed \(G\).

Choose nonzero real nonnegative functions \(f,g\in C_c^\infty(\mathbf R)\),
with \(f\) nonconstant, and define
\[
 \eta_0=f(s_1)g(s_2),\qquad
 \beta_0=f(s_1)g(s_2)\,d's_2\wedge d''s_2.
 \tag{VE24.2}
\]
These are compact original smooth objects on \(Y\): the inverse
image of a compact rectangle of torus norms is compact. The standard
original torus integration calibration gives
\[
 \begin{split}
 I:=\int_Y D\eta_0\wedge\beta_0
 &=\left(\int_{\mathbf R}f''f\,ds_1\right)
        \left(\int_{\mathbf R}g^2\,ds_2\right)\\
 &=-\left(\int_{\mathbf R}(f')^2\,ds_1\right)
        \left(\int_{\mathbf R}g^2\,ds_2\right)<0.
 \end{split}
 \tag{VE24.3}
\]
Only the first diagonal Hessian term survives its wedge with
\(\beta_0\). The signs use the same positive scalar convention
\(D=d'd''\) as VP and the frozen product calibration. In particular,
\(D\eta_0\) is a nonzero original smooth current on \(Y\).

\subsection{The polar has the full original provenance}

Choose a compact original smooth disc cutoff \(\chi_0\), equal
to one near \(T=0\), and put
\(\eta=p^*\eta_0\,\chi_0\in\mathcal A_c^{0,0}(X)\).
For formal degree-two polynomials define
\[
 \Lambda(P)=c\langle\mathcal A_Q(P),\eta\rangle,
 \qquad 0<c\le 1/\max(1,K_\eta),
 \tag{VE24.4}
\]
where \(L\Subset X\) and \(K_\eta<\infty\) are a fixed
carrier-CLN stage and constant furnished by KC4 for this actual
background and test with the full absolute carrier \(G=X\).
KC2 and KC6, or the exact canonical-polar calculation AP1--AP2,
show that \(\Lambda\) is a genuine full unit polar for
\(q^{X,2}_{G,L}\). This quantifies over every
original \(\mathcal D(X)\) function, all three families of
permitted relations, and all actual refinements. It is not a
real-tensor model declared to be a polar.

Its ordinary pure-Hessian representing current is
\[
 T_\Lambda=cD\eta\wedge Q=c\,i_*(D\eta_0).
 \tag{VE24.5}
\]
It has degree \((2,2)\), bidimension \((1,1)\), compact support,
and both separate boundaries zero. For any smooth original
\(u,v\in\mathcal D(X)\), the original signed recursion and
integration by parts give
\[
 \Lambda([u][v])
       =\langle T_\Lambda,u\,Dv\rangle.
 \tag{VE24.6}
\]
The positive sign is the frozen rule
\(\langle D(uR),\eta\rangle=\langle uR,D\eta\rangle\)
for a separately closed current \(R\), here applied to
\(R=Dv\wedge Q\). Thus the first, intrinsic Hessian-representation
hypothesis of VP is satisfied. In particular no unproved local
harmonicity property of abstract polars is being invoked.

\subsection{Shrinking original cutoffs and product calibration}

Choose one smooth real profile \(\rho\) with
\(\rho=1\) on \(( -\infty,-1]\), \(\rho=0\) on \([0,\infty)\).
For integers \(N\ge1\), set
\[
 \chi_N(T)=\rho(\log|T|+N)\quad(T\ne0),\qquad \chi_N(0)=1.
 \tag{VE24.7}
\]
Each \(\chi_N\) is an original compact smooth scalar on the
disc. It is constant on an ordinary neighborhood of zero and
uses the genuine unit \(T\) on its transition annulus. On that
annulus its curvature is
\(\rho''(\log|T|+N)d'\log|T|\wedge d''\log|T|\). Consequently
\[
 \int_{\mathbf D}D\chi_N=0,\qquad
 \|D\chi_N\|_{\rm TV}=M_\rho:=\int_{\mathbf R}|\rho''(u)|\,du.
 \tag{VE24.8}
\]
The supports of \(\chi_N\) lie in \(\{|T|\le e^{-N}\}\).
There is no requirement that these radii lie in \(|K^\times|\).

We shall use the following exact original product calculation. If
\(\tau\) is an original compact smooth top form on the base, then
\[
 \int_{Y\times\mathbf D}p^*\tau\wedge D\chi_N
         =\left(\int_Y\tau\right)\left(\int_{\mathbf D}D\chi_N\right)=0.
 \tag{VE24.9}
\]
The equality is the one-center annular fiber integration in the
frozen SP calibration, equivalently DP10 and DP13 on the compact
support of this separated form. It holds for arbitrary original
base forms, with their actual local moment refinements; the base
need not be replaced by its two-coordinate real image. The same
calculation for the nonnegative annular density \(|\rho''|\)
and a compact positive top form \(\Omega\) on the base gives
mass \(M_\rho\int_Y\Omega\). By the frozen finite coefficient
domination, any finite collection of compact smooth top forms
\(\tau_j\) has such positive controls for its coefficient
variations. Therefore there are finite constants \(M_j\),
independent of \(N\), such that
\[
 \left|\int e\,p^*\tau_j\wedge D\chi_N\right|
       \le M_j\|e\|_{\operatorname{supp}\tau_j\times\{|T|\le e^{-N}\}}.
 \tag{VE24.10}
\]
The product in the supremum denotes the corresponding closed
analytic tube. The multiplier \(e\) may be continuous. Formula
(VE24.10) uses its existing coefficient-Radon measure, not a new
product of currents. Formula (VE24.9) also holds after a fixed
continuous base multiplier: approximate it uniformly on the
compact base support by original smooth scalars and apply this
same bound. This is the exact continuous-test extension used below.

\subsection{No continuous-form primitive, even locally at the section support}

\begin{proposition}[The primitive-existence shortcut fails for a genuine polar]
Let \(y_0\in\operatorname{supp}(D\eta_0)\). On no ordinary
neighborhood \(O\) of \(i(y_0)\) is there a finite expression
\[
                 V=\sum_{j=1}^m a_j\theta_j,
 \qquad a_j\in C(O,\mathbf R),\quad
       \theta_j\in\mathcal A^{1,1}(O),
 \tag{VE24.11}
\]
whose associated ordinary current satisfies
\(D[V]=T_\Lambda|_O\) on every original compact smooth test.
In particular \(T_\Lambda\) has no primitive of the specified
compact continuous-coefficient type VP24.2 on \(X\).
\end{proposition}

\begin{proof}
Suppose (VE24.11) were given. The inverse image
\(V_0=i^{-1}(O)\) is an ordinary neighborhood of \(y_0\).
By the definition of current support there exists an original
\(\beta\in\mathcal A_c^{1,1}(V_0)\) such that
\[
                         J=\langle D\eta_0,\beta\rangle\ne0.
 \tag{VE24.12}
\]
Choose a compact original affinoid union \(A\Subset V_0\)
whose ordinary interior contains \(\operatorname{supp}\beta\).
Compactness of the actual product over \(A\), and the fact that
\(O\) contains its entire zero section, give \(r>0\) with
\(p^{-1}(A)\cap\{|T|\le r\}\subset O\).
One may first take a fixed larger closed disc and then use the
positive minimum of \(|T|\) on its compact complement of \(O\).
The zero locus of \(T\) in this actual product is exactly the
section. No topology of pairs of analytic points is used.

The exact section-horizontality lemma CC1--CC2, applied to the
finitely many \(\theta_j\), allows a smaller uniform \(r\)
for which, near this compact tube,
\[
                     \theta_j=p^*\theta_j^0,
                \qquad \theta_j^0=i^*\theta_j.
 \tag{VE24.13}
\]
For clarity, this follows from original analytic unit
presentations: every unit \(u\) occurring near the section has
unit restriction \(u_0\), and the ordinary inequality
\(|u-p^*u_0|<|p^*u_0|\) makes their norms equal.
The original forms and all their coefficient expressions thus
become horizontal. A finite compact cover gives one radius.
This is equality of actual original forms, not of a selected
real coefficient array.

For all large \(N\), the test
\(\omega_N=\chi_Np^*\beta\) is compactly supported in \(O\).
Every term in \(D\omega_N\), other than
\(D\chi_N\wedge p^*\beta\), has after wedging with
\(p^*\theta_j^0\) horizontal degree greater than two in at
least one column. For example the two first-derivative terms
have horizontal degrees \((3,2)\) and \((2,3)\), and the
\(\chi_Np^*D\beta\) term has degree \((3,3)\).
They vanish as pullbacks of actual forms on the two-dimensional
base. Thus the full differentiation, including its cutoff terms,
gives
\[
 \begin{split}
 \langle D[V],\omega_N\rangle
 &=\langle[V],D\omega_N\rangle\\
 &=\sum_j\int a_j\,p^*(\theta_j^0\wedge\beta)\wedge D\chi_N.
 \end{split}
 \tag{VE24.14}
\]
There is a positive double-transposition sign. No derivative
of a merely continuous \(a_j\) is introduced.

On the fixed compact tube the continuous function
\(a_j-p^*(i^*a_j)\) is zero on the section. Its supremum on
\(p^{-1}(A)\cap\{|T|\le e^{-N}\}\) tends to zero:
otherwise compact nested closed sets on which its absolute
value is at least some positive constant would meet the zero
section, a contradiction. This proof includes all correlated
points of the original analytic product.

Replace \(a_j\) in (VE24.14) by its horizontal continuous
section trace. The corresponding integral is exactly zero by
(VE24.9) and its continuous-multiplier extension. The error
tends to zero by (VE24.10), because the finitely many base
forms \(\theta_j^0\wedge\beta\) are fixed and compactly
supported. Hence
\[
                       \langle D[V],\omega_N\rangle\longrightarrow0.
 \tag{VE24.15}
\]
On the other hand (VE24.5) and \(\chi_N(0)=1\) give
\[
             \langle T_\Lambda,\omega_N\rangle=cJ\ne0
                    \quad\hbox{for every large }N.
 \tag{VE24.16}
\]
This contradicts the asserted all-original-test equality.
The global primitive in VP24.2 restricts to (VE24.11) near
any point after using its finitely many scalar domains, so it
is excluded as well. The point \(y_0\) exists by (VE24.3).
\end{proof}

\subsection{The same shrinking probes exclude alternative VP representatives}

The preceding contradiction concerned the explicit canonical current.
We now make each of its global detecting tests an exact Hessian test
of the same full polar. This excludes a different current with a VP
primitive without invoking any uniqueness theorem for current
representatives. Choose the already permitted profile \(\rho\)
in VE24.7 to take values in \([0,1]\), and set
\[
 u_N=\chi_Np^*\eta_0,\qquad v=s_2^2/2,
 \qquad u_NDv=\chi_Np^*\beta_0=:\omega_N.
 \tag{VE24.17}
\]
Here \(v\) is a global finite continuous smooth original psh
function; it need not have compact support. The actual test
\(u_NDv\) is compact because \(u_N\) is compact.

For completeness, \(u_N\in\mathcal D(X)\) follows from an
explicit global decomposition. For a nonnegative
\(k\in C_c^\infty(\mathbf R)\), choose
\(M\ge\tfrac12\|k''\|_\infty\). Then
\[
 k(s_j)=\bigl(k(s_j)+Ms_j^2\bigr)-Ms_j^2
 \tag{VE24.18}
\]
is a difference of nonnegative original psh functions, since
both profiles are convex in the global unit logarithm. Apply
this to \(f\) and \(g\). For the disc factor put
\[
 q_N=\log\max\{e^{-N-2},|T|\}+N+2\ge0.
 \tag{VE24.19}
\]
It is original continuous psh on the entire disc, including
\(T=0\). For
\(M\ge\tfrac12\|\rho''\|_\infty\) and
\(C\ge\|\rho'\|_\infty\), both real functions
\[
 z\longmapsto\rho(z-2)+Mz^2+Cz,
 \qquad z\longmapsto Mz^2+Cz
 \quad(z\ge0)
 \tag{VE24.20}
\]
are nonnegative, convex and nondecreasing. Their compositions
with \(q_N\) are original psh and differ by exactly \(\chi_N\).
Below the truncation threshold both cutoff descriptions equal
one; above it \(q_N-2=\log|T|+N\). The elementary identity
\[
 ab=\tfrac12\bigl((a+b)^2-a^2-b^2\bigr)
       \qquad(a,b\ge0\hbox{ original psh})
 \tag{VE24.21}
\]
expresses each product as a difference of nonnegative psh
functions. Repeated finite signed expansion, exactly as in
IL1, proves \(u_N\in\mathcal D(X)\). The displayed formula
also makes \(u_N\) original smooth. This proof does not apply
an affinoid compact-smooth-core statement to the unbounded
whole torus; every decomposition is on the unchanged \(X\).

Suppose there were any ordinary current \(T'\) and a finite
compact primitive of the specified VP24.2 type satisfying
\[
 \begin{gathered}
 T'=D[V],\qquad V=\sum_j a_j\xi_j,
 \quad \xi_j\in\mathcal A_c^{1,1}(X),\quad
 a_j\text{ continuous near }\operatorname{supp}\xi_j,\\
 \Lambda([u][w])=\langle T',uDw\rangle
       \quad\text{for all original smooth }u,w\in\mathcal D(X).
 \end{gathered}
 \tag{VE24.22}
\]
Then (VE24.17) and the canonical representation (VE24.6) give
\[
 \langle T',\omega_N\rangle
       =\Lambda([u_N][v])
       =\langle T_\Lambda,\omega_N\rangle=cI\ne0.
 \tag{VE24.23}
\]
But the already proved argument (VE24.13)--(VE24.15), using
the fixed base detector \(\beta_0\), gives
\(\langle D[V],\omega_N\rangle\to0\). If a coefficient
\(a_j\) is initially defined only near the compact support
of \(\xi_j\), multiply it by a compact smooth scalar equal
to one near that support and extend continuously by zero.
This leaves \([V]\) unchanged and permits the same compact
tube argument. Thus (VE24.22) is impossible.

This final step uses only the exact Hessian tests (VE24.17).
It does not assert that all smooth current tests are generated
by Hessian tests, or identify \(T'\) with \(T_\Lambda\).
The universal existential supply of both VP hypotheses is
therefore false for the displayed genuine full polar, even
allowing a different ordinary representing current.

\paragraph{What this does and does not decide.}
The current \(T_\Lambda\) does have the compact singular
primitive \(c\eta Q\), since \(D(c\eta Q)=T_\Lambda\).
It is the finite continuous-coefficient smooth-form class that
fails. Nothing here excludes an appropriately defined singular
primitive interface, nor proves that a different representation
of \(\Lambda\) could not be useful. We do not infer uniqueness
of arbitrary current representatives from Hessian tests.
The local obstruction applies to the explicit canonical current
(VE24.5). The final Hessian-test argument additionally excludes
any global pair of representative and primitive satisfying
(VE24.22).

The background \(Q\) is carried by a proper analytic section;
it is an allowed middle-positive background, not an arbitrary
conormal tensor. The example is not offered as an Abhyankar
model counterexample or as an actual signed PL-zero detected
by a full polar. The purpose is to delimit the sufficient
criterion VP: its two hypotheses cannot always be supplied
together, even in the action-generated situation. A future primitive interface may retain genuine
singular carrier factors such as \(Q\); no theorem producing
such a representation for every full polar is supplied here.

\paragraph{Construction and review roles.}
The polar-duality researcher constructed the section-supported
example, shrinking-test proof and exact global Hessian probes. The root checked the
mechanism, specified the exact provenance and degree guards,
and requested the full product-calibration argument. Both are
constructors. Author-separated review is recorded separately
against the exact source bytes; no constructor reviews this
new argument as independent.
